% CEE Lecture 2: k-NN lives on in retrieval-augmented generation (Lewis et al., 2020), drawn with
% the course's encoder-decoder vocabulary: the question is encoded, its k nearest passages are
% found in an index, and a language model (the decoder) answers from them.
\documentclass[tikz,border=4pt]{standalone}
\input{diagram_style}
\begin{document}
\begin{tikzpicture}
  \node[data] (q) at (0,0) {question};
  \node[op] (enc) at (3.3,0) {encoder};
  \node[data] (z) at (6.1,0) {vector $\mathbf{z}$};
  \node[focus, text width=2.6cm] (knn) at (9.6,0) {$k$ nearest\\passages};
  \node[op, text width=2.7cm] (lm) at (13.6,0) {language model\\(decoder)};
  \node[data] (a) at (17.2,0) {answer};
  \node[data, text width=2.6cm, align=center] (idx) at (9.6,-2.6) {index of\\documents};
  \draw[flow] (q) -- (enc);
  \draw[flow] (enc) -- (z);
  \draw[flow] (z) -- (knn);
  \draw[flow] (idx) -- (knn);
  \draw[flow] (knn) -- (lm);
  \draw[flow] (lm) -- (a);
  \draw[flow] (q.north) to[out=35, in=145] (lm.north);
  % the picture of "near": a query among stored vectors, its k = 3 nearest ringed
  \begin{scope}[shift={(4.1,-3.0)}]
    \foreach \x/\y in {-1.6/0.3, -1.1/-0.6, -0.5/0.75, 0.15/-0.35, 0.55/0.45, 1.2/-0.75,
        1.55/0.6, -1.7/-0.9, 0.9/0.05, -0.2/-0.95}
      \fill[muted] (\x,\y) circle (2.2pt);
    \foreach \x/\y in {0.15/-0.35, 0.55/0.45, 0.9/0.05}
      \draw[accent, line width=1.1pt] (\x,\y) circle (6pt);
    \fill[alertc] (0.45,0.0) circle (3.6pt);
    \draw[accent, line width=1pt, densely dotted] (0.45,0.0) circle (0.62);
  \end{scope}
  \draw[leader] (5.3,-2.5) -- (z.south);
  \node[note, anchor=north west, text width=6.6cm] at (11.6,-1.55)
    {the encoder decides what \emph{near} means, as the groups did for concrete;\\[2pt]
     with no truly similar passage, the answer can still sound sure};
\end{tikzpicture}
\end{document}
